\subsubsection*{Ordinary Milnor elements and infinitely divisible residues}

The following intermediate result uses the refined weight bound in
\citet[Proposition~3.7, p.~16]{BCTSchreiederDivisibility2026} to extend the
conclusion of \citet[Lemma~4.1, pp.~11--12]{BSchreiederMilnorTate2026} beyond
its stated degree and unramified-vanishing hypotheses.  No novelty claim
is made, and this is not a resolution of the invariant-colimit gate.

\begin{proposition}[Kernel of the ordinary generic Milnor regulator]
\label{prop:B-Milnor-divisible-kernel}
\status{Proved intermediate statement; the universal B.1 gate is unchanged}
Let $k=\overline{\mathbf F}_p$, let $K/k$ be a finitely generated field
extension, and let $\ell\ne p$.  Define ordinary generic continuous
cohomology by
\[
 H^n_\eta(K,\mathbf Z_\ell(n))
   =\varinjlim_{\varnothing\ne U\subset X}
      H^n_{\mathrm{cont}}(U,\mathbf Z_\ell(n)),
\]
where $X$ is a model of $K$ and the limit can be taken over its smooth
dense opens.  For the ordinary cycle map
\[
 \rho_n:K_n^M(K)\longrightarrow
                 H^n_\eta(K,\mathbf Z_\ell(n)),
\]
one has, for every $n\ge0$,
\[
 \ker\rho_n=\bigcap_{r\ge1}\ell^rK_n^M(K).
 \tag{B.M1}\label{eq:B-Milnor-kernel}
\]
For $n\ge1$, suppose moreover that $X/k$ is smooth projective and that
$\alpha\in K_n^M(k(X))$ has infinitely $\ell$-divisible residue at every
prime divisor of $X$.  Then $\alpha$ is infinitely $\ell$-divisible.
The second assertion deliberately excludes $n=0$.
\end{proposition}

\begin{proof}
The cycle map is defined on finite-type opens before passing to the
colimit.  It can be constructed from products of Kummer classes of
symbols, or from the compatible motivic cycle map in
\citet[\S6.1, equations~(6.2)--(6.4)]{BCTSchreiederDivisibility2026}.
It commutes with residues, finite-extension restriction and transfer,
and reduction modulo $\ell^r$.  Reduction on the field is the
norm-residue isomorphism.  Consequently $\rho_n(\alpha)=0$ implies
$\alpha\in\ell^rK_n^M(K)$ for every $r$.

We prove the reverse implication by induction on $n$, simultaneously
for all finitely generated extensions of $k$.  The case $n=0$ is the
injection $\mathbf Z\hookrightarrow\mathbf Z_\ell$.
Assume the assertion in degree $n-1$.  First suppose that $K$ has a
smooth projective model $X/k$, and let $\alpha$ be infinitely
$\ell$-divisible.  Every residue $\partial_D\alpha$ is infinitely
$\ell$-divisible in $K_{n-1}^M(k(D))$.  The induction hypothesis for
the residue field $k(D)$ therefore gives
\[
 \partial_D\rho_n(\alpha)
       =\rho_{n-1}(\partial_D\alpha)=0
\]
in ordinary generic continuous cohomology for every prime divisor $D$.

Localization and codimension-one purity place the rationalized class
in the ordinary refined unramified group
\[
 H^n_{0,\mathrm{nr}}(X,\mathbf Q_\ell(n))
 =\operatorname{im}\bigl(
 H^n(F_1X,\mathbf Q_\ell(n))\longrightarrow
 H^n(F_0X,\mathbf Q_\ell(n))\bigr).
\]
Here is the finite-support detail.  Start with one open on which the
class is defined.  It misses finitely many prime divisors.  The zero
generic residue on each of them vanishes already on a dense open of
that divisor.  Delete the finitely many remaining closed subsets of
codimension at least two in $X$, including singular and intersection
loci.  The localization sequence then lifts the class across all
prime divisors.  This argument involves an ordinary filtered
colimit; it makes no assertion about commuting that colimit with an
inverse limit.  The precise refined-unramified definition is
\citet[equations~(2.7)--(2.8), p.~11]{BCTSchreiederDivisibility2026}.

By the cited Proposition~3.7, the weights of this rational group lie
in
\[
 [-n,\max(-n,-2)],
\]
which is strictly negative for $n\ge1$.  Choose a finite field over
which $X$ and $\alpha$ are defined.  The ordinary Milnor element
$\alpha$ has a finite Frobenius orbit, as does its regulator.  The
finite-dimensional span of that orbit has only roots of unity as
Frobenius eigenvalues and thus cannot be a nonzero submodule of this
strictly negative-weight group.  Hence $\rho_n(\alpha)$ is zero
rationally.  The ordinary integral generic continuous group is
torsion-free by \citet[Theorem~2.2 and Corollary~2.3,
pp.~6--7]{BSchreiederMilnorTate2026}; the change of twist used for
torsion-freeness is allowed because $k$ contains all $\ell$-power
roots of unity.  Thus $\rho_n(\alpha)=0$ integrally.

For a general $K$, choose a projective model over a finite field and
a smooth projective alteration
\citep[Theorem~4.1 and Remark~4.2]{c_dejong1996}.
After choosing a geometric component and, if necessary, enlarging
the finite field of definition, its function field $L$ is a finite
extension of $K$.  The pullback of $\alpha$ is infinitely
$\ell$-divisible and its regulator vanishes by the just-proved
smooth-projective case.  Restriction followed by corestriction is
multiplication by $[L:K]$ on ordinary generic continuous cohomology:
this is finite-\'etale transfer on suitable dense opens for the
separable part, and the purely inseparable part is handled by
universal-homeomorphism invariance and its degree.  Therefore
\[
 [L:K]\rho_n(\alpha)=0.
\]
Torsion-freeness gives $\rho_n(\alpha)=0$.  This proves the induction
statement for arbitrary residue fields as well as for $K$ itself.
No degree-prime-to-$\ell$ condition is needed in this argument.

Finally, take an ordinary $\alpha$ on a smooth projective $X$ whose
residues are infinitely $\ell$-divisible, without initially assuming
the same of $\alpha$.  The proved degree-$(n-1)$ assertion again
kills every residue of its regulator.  The same localization,
negative-weight and finite-orbit argument gives
$\rho_n(\alpha)=0$.  Formula~\eqref{eq:B-Milnor-kernel} then gives
infinite $\ell$-divisibility.  This last argument requires $n\ge1$:
in degree zero all residues vanish whereas $K_0^M=\mathbf Z$ is not
divisible.
\end{proof}

\begin{corollary}[The separated Milnor quotient is free]
\label{cor:B-Milnor-free-quotient}
\status{Proved intermediate structural consequence; no new universal gate}
In the notation of Proposition~\ref{prop:B-Milnor-divisible-kernel}, set
\[
 D_\ell(K_n)=\bigcap_{r\ge1}\ell^rK_n^M(K),\qquad
 M_n(K)=K_n^M(K)/D_\ell(K_n).
\]
Then $D_\ell(K_n)$ is $\ell$-divisible and $M_n(K)$ is a free abelian
group of at most countable rank.  Moreover, $D_\ell(K_n)$ is independent
of the prime $\ell\ne p$.  Thus one canonical quotient $M_n(K)$
induces the canonical isomorphisms of completions
\[
 \widehat{K_n^M(K)}_\ell\simeq\widehat{M_n(K)}_\ell
 \qquad(\ell\ne p),
\]
equivariant for the Frobenius action from any finite-field model.
No saturated residue embedding or equivariant permutation basis is
asserted, and a common integral module does not by itself make its
completed Frobenius spectra independent of the prime.
\end{corollary}

\begin{proof}
Since $D_\ell(K_n)=\ker\rho_n$ and the regulator target is torsion-free,
$M_n(K)$ is torsion-free.  If $a\in D_\ell(K_n)$, choose $b$ with
$\ell b=a$.  Then $\ell\rho_n(b)=0$, so $\rho_n(b)=0$ and
$b\in D_\ell(K_n)$.  Thus $D_\ell(K_n)$ is $\ell$-divisible.

We prove freeness by induction on $n$ for all finitely generated $K/k$.
The case $n=0$ is $M_0(K)=\mathbf Z$.  For arbitrary $K$, take a finite
extension $L/K$ with a smooth projective model, supplied by an
alteration as in the proposition.  Restriction induces an injection
$M_n(K)\hookrightarrow M_n(L)$: if the restricted class has zero
regulator, transfer kills $[L:K]\rho_n(a)$, and torsion-freeness gives
$\rho_n(a)=0$.  It therefore suffices to prove freeness for the
function field of a smooth projective $X$.

For $n\ge1$, the residue assertion of the proposition gives an injection
\[
 M_n(k(X))\lhook\joinrel\longrightarrow
       \bigoplus_{D\in X^{(1)}}M_{n-1}(k(D)).
\]
Each ordinary Milnor element has only finitely many nonzero residues,
so this is a direct sum.  By the induction hypothesis its summands
are free abelian; hence so is the direct sum, and every subgroup of
a free abelian group is free.  All the fields and indexing sets here
are countable because $\overline{\mathbf F}_p$ is countable.  This
completes the induction.

Finally, divisibility of $D_\ell(K_n)$ gives, for every $r$, the
canonical equivariant isomorphism
\[
 K_n^M(K)/\ell^rK_n^M(K)\simeq M_n(K)/\ell^rM_n(K).
\]
Taking inverse limits proves the completion assertion.  None of the
argument makes a residue embedding $\ell$-saturated or chooses a
Frobenius-stable basis.

For independence of the kernel, let $\ell,\ell'\ne p$.  The image of
$D_\ell(K_n)$ in the free abelian group
$K_n^M(K)/D_{\ell'}(K_n)$ is $\ell$-divisible.  Such an image is zero:
an element of a free abelian group has finitely many integer
coordinates and cannot be divisible by every power of $\ell$ unless
it vanishes.  Thus $D_\ell(K_n)\subset D_{\ell'}(K_n)$, and symmetry
gives equality.  The common subgroup is divisible by every prime
different from $p$, and its common quotient gives all the displayed
completion isomorphisms.
\end{proof}

\paragraph{The remaining obstruction.}
The result concerns ordinary Milnor elements, whose Frobenius
orbits are finite.  It does not identify all continuous generic
classes, or all completed Milnor classes, with such elements.
In the first B.1 gate, a residue of
$H^4(F_0\mathbf P^4_k,\mathbf Q_\ell(3))$ lies in degree three
and twist two on a boundary threefold.  Proposition~3.7 gives the
unramified weight interval $[-1,0]$ there, so weight zero is not
excluded.  Replacing that twist by the diagonal twist three would
change the problem.  Moreover, Proposition~\ref{prop:B-limit-diagnostics}
shows that completion of a locally finite Frobenius module can
create non-root-of-unity eigenvectors; its abstract function module
is not asserted to be a Milnor module.

The additional weight input in the proof is essential to this
argument.  Inverse-limit divisibility statements alone do not
correct residues.  For example, the complex
\[
 C^0=\mathbf Z\ \hookrightarrow C^1=\mathbf Z[1/\ell],\qquad C^2=0
\]
has $H^0(C)=0$, divisible $H^1(C)=\mathbf Z[1/\ell]/\mathbf Z$,
zero primary torsion in its terms, and $H^1(C/\ell^r)=0$ for every
$r$.  Nevertheless, $1\in C^0$ is not infinitely $\ell$-divisible
while its image in $C^1$ is.  This is a formal obstruction to that
inference, not a geometric counterexample.
